\documentclass[twocolumn,prd,showpacs,superscriptaddress]{revtex4-2}
\usepackage{amsmath,amssymb,amsthm}
\usepackage{graphicx}
\usepackage{hyperref}
\usepackage{booktabs}
\usepackage{multirow}
\usepackage{tabularx}
\usepackage{float}
\usepackage{color}
\usepackage{xcolor}
% Graphics path for figures
\graphicspath{{figures/}}

% Theorem/lemma environments
\newtheorem{proposition}{Proposition}
\newtheorem{lemma}{Lemma}

% Commutator notation
\newcommand{\comm}[2]{[#1, #2]}
\newcommand{\norm}[1]{\lVert #1 \rVert}
\newcommand{\inner}[2]{\langle #1, #2 \rangle}
\newcommand{\R}{\mathbb{R}}
\newcommand{\N}{\mathbb{N}}
\newcommand{\E}{\mathbb{E}}
\newcommand{\Var}{\operatorname{Var}}
\newcommand{\Tr}{\operatorname{Tr}}
\newcommand{\diag}{\operatorname{diag}}
\newcommand{\rank}{\operatorname{rank}}
\newcommand{\sgn}{\operatorname{sgn}}

% Greek letters
\newcommand{\rhohat}{\hat{\rho}}
\newcommand{\Pihat}{\hat{\Pi}}
\newcommand{\Gammahat}{\hat{\Gamma}}
\newcommand{\thetahat}{\hat{\theta}}

% Misc
\newcommand{\Msun}{M_\odot}
\newcommand{\Mchirp}{\mathcal{M}_c}
\newcommand{\chirp}{\mathcal{M}}
\newcommand{\Lambdatilde}{\tilde{\Lambda}}
\newcommand{\chieff}{\chi_{\text{eff}}}
\newcommand{\Dl}{D_L}
\newcommand{\fisco}{f_{\text{ISCO}}}
\newcommand{\fmin}{f_{\text{min}}}
\newcommand{\fmax}{f_{\text{max}}}
\newcommand{\fbif}{f_{\text{bif}}}
\newcommand{\vd}{v_{d}}          % degeneracy direction
\newcommand{\rhosq}{\rho^{2}}    % discriminating SNR^2

\newcommand{\figref}[1]{Fig.~\ref{fig:#1}}
\newcommand{\tabref}[1]{Table~\ref{tab:#1}}
\newcommand{\secref}[1]{Sec.~\ref{sec:#1}}
\newcommand{\eqnref}[1]{Eq.~(\ref{eq:#1})}

% Manifest placeholder macro — replaced by generated values from analysis pipeline
\newcommand{\manifest}[1]{\textbf{[#1]}}


\begin{document}

\title{The ($\rho$, $\Pi$) Bridge to Gravitational-Wave Inference: \\
Commutator-Modulated Dynamics and the Geometry of Sequential Bayesian Inference}

\author{Skye Laflamme}
\author{Axioma}
\affiliation{RavenNest Science, Santa Cruz, CA, USA}

\date{June 2026}

\begin{abstract}
We present the ($\rho$, $\Pi$) bridge: a formal connection between the
commutator-modulated dynamical system from the S3 agent framework and
the geometry of sequential Bayesian inference in gravitational-wave
astronomy.  The central claim is that the ($\rho$, $\Pi$) dynamics
describe the evolution of the posterior distribution over binary
orbital parameters as data accumulates across frequency, not the
orbital dynamics of the binary itself.  We distinguish two logically
independent objects: the exact posterior trajectory
$\{\rho_f\}_{f_{\text{min}}\le f\le f_{\text{max}}}$, whose geometric
properties are observer-independent features of the inference problem;
and the deterministic bias of single-pass sequential-Gaussian
(assumed-density filtering) inference relative to batch sampling, which
our formalism characterizes quantitatively.

We present three lines of evidence.  \textbf{Static evidence:} The
GW170817 posterior in the ($q$, $\chi_{\text{eff}}$) plane exhibits an
extended degeneracy ridge, not a bifurcation.  The displacement between
the prior-conditioned posterior means aligns with the Fisher
degeneracy eigenvector, consistent with the prior-lensing prediction
of the formalism.  The alignment is assessed against the null
distribution of a random direction in the four-dimensional
$\mathcal{M}_c$-orthogonal subspace, yielding a $p$-value
$p = \manifest{T3: alignment p-value}$
--- a falsifiable statement that can fail for future events.
\textbf{Dynamic evidence:} A Fisher information sweep across 18
frequency cutoffs from 20--500~Hz reveals a
$\manifest{T6: C(f) growth factor}$ growth in the commutator
misalignment $C(f)$ across the chirp window
($126$~s $\rightarrow$ $0.03$~s before merger).  The growth law is
derived analytically from post-Newtonian power counting; the
predicted slope matches the numerical sweep to within
$\manifest{T7: slope agreement band}$.
\textbf{Forward model:} A derived weight trajectory, slaved to the
exact accumulated likelihood ratio between the ridge endpoints, predicts
a resolvable peak-drop signature in time-resolved inference at
$f_* \approx \manifest{T9: predicted f*}$ Hz, testable in an
in-software ordering experiment.

The systematic offset formula
$\Delta\theta_{\text{sys}} \approx \Gamma(f_{\text{max}})^{-1}
\int [\Gamma(f'), \Pi(f')]\,df' \cdot v_{\perp}$
is reframed as the bias of single-pass sequential-Gaussian filtering
relative to exact batch inference, with the commutator governing how
approximation error leaks into the degeneracy direction.  A dual-PSD
Fisher sweep shows the commutator integral is PSD-insensitive at the
$\sim\!1\%$ level (ratio $\manifest{T8: dual-PSD ratio}$),
confirming that the bias formula's inputs are detector-universal.
The decisive test---an in-software comparison of batch, sequential-forward,
and permuted-order inference on a single injection---is described;
it replaces the previously proposed multi-detector comparison, which
is shown to be unperformable given the $\sim\!1\%$ effect size and
coherent-network PE design.
\end{abstract}

\maketitle

\section{Introduction}
\label{sec:intro}

The analysis of gravitational-wave data from compact binary
coalescences~\cite{maggiore2007gw} is a problem of sequential Bayesian
inference.  As a signal sweeps across the detector's sensitive band
from low frequencies to merger, the likelihood surface over the
binary's orbital parameters evolves: initially broad and degenerate,
then progressively sharper as higher post-Newtonian (PN) orders become
resolvable.  This evolution is not merely a technical detail---it is a
dynamical process with its own geometry, its own timescales, and
potentially its own observable consequences for the interpretation of
parameter estimates.

In this paper, we develop a formal framework for describing this
evolution.  We define an S3 agent as a pair ($\rho$, $\Pi$), where
$\rho$ is the posterior density and $\Pi$ is the prior-plus-likelihood
structure.  The central object of this system is the commutator
$[\Gamma, \Pi]$, which measures the cross-coupling between the
best-constrained direction (the leading Fisher eigenvector, dominated
by chirp mass) and the remaining parameter subspace.  Crucially, we
distinguish two objects that are conflated in the earlier draft of
this work: (i) the \emph{exact posterior trajectory}
$\{\rho_f\}$, a one-parameter family whose geometric properties
(eigenframe rotation, commutator growth) are exact, observer-independent
features of the inference problem; and (ii) the \emph{deterministic
bias of single-pass sequential-Gaussian (assumed-density filtering)
inference} relative to batch sampling, which arises because each
Laplace projection discards non-Gaussian structure that would otherwise
interact with subsequent updates.  The commutator governs the latter
via an eigenframe-leakage mechanism that we derive in the Appendix.

The paper is organized as follows.
Section~\ref{sec:formalism} presents the formalism, the exactness
theorem, the coordinate convention, the three-way typing of Fisher
validity, and the degeneracy-ridge and prior-lensing structure.
Section~\ref{sec:static} presents the static test: ridge geometry and
prior-displacement alignment in GW170817, with a $p$-value against
the random-direction null.
Section~\ref{sec:dynamic} presents the dynamic test: the Fisher
information sweep, commutator growth, the derived power-counting
growth law, and PSD-insensitivity.
Section~\ref{sec:forward} presents the forward model for the
time-resolved signature, including the derived weight trajectory
and the non-Gaussianity index.
Section~\ref{sec:open} discusses the decisive ordering experiment
and other open problems.
Section~\ref{sec:interpretation} interprets the bridge.
An appendix provides the eigenframe-leakage derivation and
conditioning diagnostics.

\section{Formalism: Exact Trajectories and Approximate Filters}
\label{sec:formalism}

\subsection{The S3 agent pair}

We define an S3 agent as a pair ($\rho$, $\Pi$), where $\rho$ is a
time-evolving probability measure over a parameter space $\Theta$ and
$\Pi$ is the structure that updates it: the prior together with the
likelihood of the data accumulated so far.  For gravitational-wave
inference on a compact binary coalescence,
$\Theta = \{\Mchirp, q, \chieff, \Lambdatilde, \Dl\}$, and the natural
ordering variable is the frequency cutoff $f$, monotonic in time as the
signal chirps across the detection band.  We write $\rho_f$ for the
posterior given data below cutoff $f$.

\subsection{Specialisation to GW inference}

For stationary Gaussian noise with one-sided power spectral density
$S_n(f)$, the log-likelihood at cutoff $f$ is
\begin{equation}
\ln \mathcal{L}(\theta \mid d, f) = -\frac12 \int_{\fmin}^{f}
\frac{|d(f') - h(f'; \theta)|^2}{S_n(f')}\,df' + \text{const},
\label{eq:likelihood}
\end{equation}
and the Fisher information
matrix~\cite{cutler1994fisher,vallisneri2008fisher} is
\begin{equation}
\Gamma_{ij}(f) = \int_{\fmin}^{f}
\frac{\partial_i h(f';\theta)\,\partial_j h(f';\theta)^{*} + \text{c.c.}}
{2\,S_n(f')}\,df',
\label{eq:fisher}
\end{equation}
evaluated at a stated expansion point (\secref{validity}).  The
increment
\begin{equation}
\frac{d\Gamma_{ij}}{df} =
\frac{\partial_i h\,\partial_j h^{*} + \text{c.c.}}{2\,S_n(f)}
\label{eq:fisher_flow}
\end{equation}
is a deterministic flow on the cone of positive-semidefinite matrices,
driven entirely by the waveform model and the noise weighting.

\subsection{Exactness, and the two objects of study}
\label{sec:exactness}

Before introducing dynamics, we fix what can and cannot evolve.  For
stationary Gaussian noise the likelihood factorises over disjoint
frequency bands, so Bayesian updating over those bands is a product of
commuting factors:

\begin{proposition}[Order independence of exact inference]
\label{prop:exact}
Let $\{B_k\}_{k=1}^{N}$ be a partition of $[\fmin,\fmax]$ into disjoint
bands with likelihood factors $\mathcal{L}_k(\theta)$.  Then for any
permutation $\sigma$,
$\pi(\theta)\prod_k \mathcal{L}_{\sigma(k)}(\theta)
= \pi(\theta)\prod_k \mathcal{L}_{k}(\theta)$:
the exact posterior at $\fmax$ is independent of the order in which the
bands are assimilated.
\end{proposition}

Proposition~\ref{prop:exact} forbids any claim that the \emph{exact}
posterior depends on the path of information accumulation.  Two
well-defined objects survive it, and they are the two subjects of this
paper.

\paragraph{Object 1: the exact posterior trajectory.}
The one-parameter family $\{\rho_f : \fmin \le f \le \fmax\}$, where
each $\rho_f$ is the exact posterior given data below $f$, is
well defined, and the frequency ordering of this family is physically
privileged: it is the time ordering experienced by any real-time
observer of a chirping binary, since data above $f$ does not yet exist
when the signal instantaneous frequency is $f$.  The geometry of this
family---how the local information metric reorients as $f$ grows
(\secref{commutator})---is an exact, observer-independent property of
the inference problem.  Every endpoint statement is anchored by
Proposition~\ref{prop:exact}: $\rho_{\fmax}$ equals the batch
posterior.

\paragraph{Object 2: the single-pass Gaussian filter and its bias.}
Practical low-latency analyses cannot afford batch sampling and instead
propagate a Gaussian summary.  We formalise the scheme as
assumed-density filtering (ADF)~\cite{minka2001ep,opper1998online} with
a Laplace projection: maintain
$q_k = \mathcal{N}(m_k, P_k)$; on receipt of band $B_{k+1}$, linearise
$h$ about $m_k$, perform the conjugate (Gauss--Newton) update, and
re-Gaussianise.  Each projection discards non-Gaussian structure that
would otherwise interact with subsequent bands, so---unlike the exact
update---\emph{the single-pass filter is order dependent}, and its
terminal estimate carries a deterministic bias relative to the batch
posterior.  Characterising that bias, and showing that it is governed by
the commutator structure of \eqnref{commutator_def}, is the second
subject of this paper (\secref{dynamic}, Appendix~\ref{sec:appendix}).
We emphasise the division of labour: nothing in Object~1 is an
approximation, and nothing in Object~2 is a property of exact
inference.

\subsection{Coordinates and invariance}
\label{sec:coordinates}

Angles, eigenvectors, and Frobenius norms in parameter space are
meaningful only relative to a metric, and the raw parameters carry
heterogeneous units: a Euclidean angle computed with $\Mchirp$ in solar
masses changes if $\Mchirp$ is expressed in seconds.  We therefore fix,
once, the dimensionless analysis coordinates
\begin{equation}
x = \bigl(\ln \Mchirp,\; q,\; \chieff,\; \Lambdatilde/100,\; \ln \Dl\bigr),
\label{eq:coords}
\end{equation}
in which all eigendecompositions, projectors, and matrix norms in this
paper are computed; Fisher matrices and displacement vectors transform
by the Jacobian of \eqnref{coords} in the standard way.  Robustness of
every headline quantity to two alternative conventions (a
symmetric-mass-ratio variant and a prior-range whitening) is reported
alongside the primary values (\secref{dynamic}).

Three classes of quantity require no convention at all, because the
Fisher matrix is the pullback of the waveform-space inner product
$\inner{a}{b} = 4\,\mathrm{Re}\!\int a b^{*}/S_n\,df$ onto parameter
space: (i) likelihood scalars such as
$\Delta\ln\mathcal{L} = \tfrac12\,\Delta x^{\top}\Gamma\,\Delta x$;
(ii) the exact mismatch
$\inner{\delta h}{\delta h}$ between two waveforms; and (iii) the
cumulative discriminating signal-to-noise between two parameter points,
\begin{equation}
\rhosq(f) \;\equiv\; \Delta x^{\top}\,\Gamma(f)\,\Delta x
\;\simeq\; \inner{\,h(\theta_1)-h(\theta_2)\,}{\,h(\theta_1)-h(\theta_2)\,}_{<f},
\label{eq:rhosq}
\end{equation}
where the right-hand side is the exact (convention-free) quantity and
the quadratic form its second-order approximation; we report both and
use the exact form wherever the displacement is not small
(\secref{validity}).  Wherever possible, predictions in this paper are
stated in terms of these invariants.

\subsection{The projector, the commutator, and the misalignment $C(f)$}
\label{sec:commutator}

Let $\lambda_1(f) \ge \cdots \ge \lambda_5(f)$ and $v_1(f), \ldots,
v_5(f)$ be the eigenvalues and unit eigenvectors of $\Gamma(f)$ in the
coordinates of \eqnref{coords}.  Define the projector onto the
best-constrained direction,
\begin{equation}
\Pi(f) = v_1(f)\,v_1(f)^{\top},
\label{eq:projector}
\end{equation}
and the commutator
\begin{equation}
\comm{\Gamma}{\Pi} \equiv \Gamma\Pi - \Pi\Gamma
= \lambda_1 v_1 v_1^{\top} - v_1 v_1^{\top}\Gamma .
\label{eq:commutator_def}
\end{equation}
The commutator vanishes iff $\Gamma$ has no off-diagonal structure
connecting $v_1$ to its orthogonal complement at that cutoff; it is the
infinitesimal generator of the eigenframe rotation of the cumulative
Fisher matrix.  We quantify the rotation by the fractional misalignment
\begin{equation}
C(f) = \frac{\norm{\Gamma(f) - \lambda_1(f)\,\Pi(f)}_F}
{\norm{\Gamma(f)}_F},
\label{eq:cross_coupling}
\end{equation}
computed in the declared coordinates.  We further denote by
$\vd(f) \equiv v_5(f)$ the \emph{degeneracy direction}: the eigenvector
of the smallest eigenvalue.  For the binary-neutron-star systems
considered here, $\vd$ lies almost entirely in the
$(q, \chieff, \Lambdatilde)$ block at all cutoffs
[\manifest{T6: max component of $\vd$ outside the block}], reflecting
the well-known mass-ratio--spin--tidal degeneracy~\cite{chatziioannou2020bns}.
(Earlier drafts and some of the literature label this direction $v_2$
from analyses restricted to a low-dimensional subspace; we use $\vd$
throughout to avoid ambiguity about the ambient dimension.)

\subsection{The degeneracy ridge and prior lensing}
\label{sec:ridge}

At low cutoffs the likelihood \eqnref{likelihood} is nearly constant
along $\vd$: extended families of $(q, \chieff, \Lambdatilde)$
combinations produce waveforms whose accumulated discriminating
signal-to-noise \eqnref{rhosq} is below unity.  The posterior in this
subspace is therefore a \emph{degeneracy ridge}---a continuous,
elongated, prior-bounded structure---rather than a collection of
isolated modes.  We make no claim in this paper that the ridge fragments
into distinct maxima at any cutoff; whether it does is an open question
for time-resolved analysis (\secref{open}), and we deliberately avoid
bifurcation vocabulary until that analysis exists.

The ridge has an immediate, testable consequence: it determines how the
posterior responds to changes of prior.

\begin{lemma}[Prior lensing]
\label{lem:lensing}
Let the posterior be $p \propto \mathcal{L}\,\pi$ with likelihood
precision $\Gamma$ at the mean, and perturb the log-prior by a smooth
$\delta(\theta)$.  To leading order in the Laplace expansion the
posterior mean shifts by
\begin{equation}
\Delta x_{*} \approx \Gamma^{-1}\,\nabla\delta
= \sum_{i=1}^{5} \lambda_i^{-1}\,(v_i^{\top}\nabla\delta)\,v_i ,
\label{eq:lensing}
\end{equation}
which for any perturbation with generic overlap
($v_d^{\top}\nabla\delta \neq 0$) is dominated by the degeneracy
direction: $\Delta x_{*} \parallel \vd$ up to corrections of order
$\lambda_d/\lambda_{i\neq d}$.
\end{lemma}

The inverse Fisher matrix acts as a lens that focuses arbitrary prior
perturbations into the degenerate subspace.  For a hard truncation of
the prior (the case realised by the two published GW170817 spin priors)
the smooth-perturbation magnitude in \eqnref{eq:lensing} is not
controlled---the mean shift saturates at the prior boundary---but the
\emph{direction} of the shift remains that of the ridge, since
truncation removes probability mass from one end of an elongated
structure.  Lemma~\ref{lem:lensing} therefore yields a directional
prediction, tested in \secref{static}: the displacement between the
posterior means of two prior-conditioned analyses of the same data
aligns with $\vd$, with alignment assessed against the null
distribution of a random direction in the relevant subspace.

\subsection{Domain of validity of the Fisher description}
\label{sec:validity}

The matrix $\Gamma(f)$ appears in this paper in three logically
distinct roles with different validity domains.  Conflating them is the
standard failure mode of Fisher-based
arguments~\cite{vallisneri2008fisher}, so we type every use explicitly
and adhere to the typing throughout.

\paragraph{U1 --- Local information geometry (exact).}
As the pullback of the waveform-space metric, $\Gamma(f)$ and its
eigenstructure are exact local properties of the likelihood at the
expansion point, independent of the posterior's global shape.  The
misalignment $C(f)$, the commutator \eqnref{commutator_def}, and their
evolution with cutoff are U1 statements and require no Gaussianity
assumption.  Because local geometry can vary along an extended ridge,
all U1 quantities are computed at three expansion points spanning the
empirically probed ridge segment, and the spread is reported as a
systematic band [\manifest{T6: ridge-uniformity band}].

\paragraph{U2 --- Global covariance proxy (approximate; its failure is
quantified, not assumed away).}
The Laplace approximation
$\rho_f \approx \mathcal{N}(x_{*}(f), \Gamma(f)^{-1})$ is invalid along
$\vd$, where the true posterior is a prior-bounded ridge.  We do not
rely on it there; instead we \emph{measure} its failure with the
non-Gaussianity index
\begin{equation}
N(f) \;=\;
\frac{\Var_{\rho_f}\!\bigl(\vd^{\top} x\bigr)}
{\bigl(\Gamma(f)^{-1}\bigr)_{\vd\vd}} ,
\label{eq:ngauss}
\end{equation}
the ratio of the true posterior variance along the degeneracy direction
to the Fisher prediction; $N \approx 1$ certifies the Gaussian regime
and $N \gg 1$ quantifies the breakdown.  The single-pass filter bias of
\secref{exactness} accrues precisely where $N$ is large, so U2 failure
is not a caveat to this work but part of its subject matter.

\paragraph{U3 --- Pairwise discrimination (quadratic order; exact
upgrade available).}
The scalar $\rhosq(f)$ of \eqnref{rhosq} approximates the discriminating
signal-to-noise between two specified points to second order in the
displacement.  Where the displacement is not small---as for the ridge
endpoints of GW170817---we compute the exact mismatch directly from
two waveform evaluations and quote the quadratic form only alongside
its exact value [\manifest{T5: quadratic-to-exact ratio}].

Finally, we verify the typing empirically with the consistency
criterion of Ref.~\cite{vallisneri2008fisher}: at each cutoff and along
each eigendirection $v_i$, displace by one Fisher standard deviation
$\lambda_i^{-1/2}$ and compare the exact mismatch to its quadratic
prediction.  The resulting pass/fail map over $(f, v_i)$
(\figref{validity}) shows the linear-signal regime holding along $v_1$
throughout the band and failing along $\vd$ at low cutoffs, with the
onset of failure tracking $N(f)$---a direct empirical delineation of
where each role of $\Gamma$ is licensed.  Numerical conditioning of the
eigenproblem (the smallest spectral gaps, and machine-precision
derivatives in place of finite differences) is documented in
Appendix~\ref{sec:appendix}.


\section{Static Test: Degeneracy Geometry and Prior Sensitivity in GW170817}
\label{sec:static}

\subsection{The prediction}

If the posterior is shaped by a degeneracy ridge along $\vd$, then the
displacement between two prior-conditioned analyses of the same data
should align with $\vd$.  This is a nontrivial directional prediction:
the $\mathcal{M}_c$-orthogonal complement is four-dimensional, so
alignment with the \emph{specific} direction $\vd$ within that subspace
is a falsifiable claim.  The angle $\alpha$ between $\Delta x$ (the
restricted displacement in the $\mathcal{M}_c$-orthogonal subspace) and
$\vd(f_{\text{max}})$ can be expressed via the invariant cosine
\begin{equation}
\cos^2\alpha = \frac{(\vd^{\top}\Delta x)^2}{\|\Delta x\|^2},
\label{eq:cos2alpha}
\end{equation}
and under the null hypothesis that $\Delta x$ is uniformly random on
the sphere in four dimensions, $\cos^2\alpha$ follows a
$\mathrm{Beta}(1/2,\, 3/2)$ distribution.  The $p$-value
$P(\cos^2\alpha \ge \text{observed} \mid \text{random direction})$
quantifies the statistical significance of the alignment.

We additionally test whether the ridge orientation itself---the
principal axis of the posterior sample covariance within each
single-analysis posterior---matches $\vd$, and whether the contrasting
GW150914, whose Fisher matrix has a much smaller condition number in
the ($q, \chieff$) block (no tidal sector, shorter inspiral), shows a
correspondingly smaller prior-induced displacement.

\subsection{Data and method}

The LIGO open science center (GWOSC) provides posterior samples for
all confident gravitational-wave detections.  For
GW170817~\cite{abbott2017gw170817,abbott2019gw170817_params}, two
separate posterior analyses were released with different spin priors:
lowSpin ($\chi < 0.05$) and highSpin ($\chi < 0.89$).  Both used the
same strain data (Hanford $+$ Livingston, GPS times
1187008882--1187008890) with the IMRPhenomPv2NRT waveform approximant;
a TaylorF2 $+$ NRTidal release provides a cross-approximant check.
Neither posterior is bimodal in $(q, \chieff)$; each is a continuous,
banana-shaped degeneracy ridge.

We extract posterior samples in the declared coordinates
(\eqnref{coords}), compute the sample covariance eigenvectors (for
ridge orientation), and the displacement vector $\Delta x = \mu_{\text{highSpin}} - \mu_{\text{lowSpin}}$.
All quantities are reported with bootstrap 5--95\% confidence
intervals from $N_{\text{boot}} = 1000$ resamples.

\subsection{Ridge orientation (Test A)}

For each prior analysis separately, we compute the principal axis of
the sample covariance in the $(q, \chieff, \Lambdatilde)$ block and
compare it to $\vd(f_{\text{max}})$.  The alignment angle is
$\manifest{T2: ridge alignment angle}$ for the lowSpin analysis and
$\manifest{T2: ridge alignment angle (highSpin)}$ for the highSpin
analysis.  The pre-registered falsification threshold is
misalignment $> 20^\circ$ at 95\% confidence; it is not exceeded.

\subsection{Prior-displacement alignment (Test B)}

The displacement vector $\Delta x$ has components:
\begin{itemize}
  \item $\Delta \ln \Mchirp = \manifest{T3: delta ln Mc}$ (consistent
  with zero to within bootstrap uncertainty, confirming the
  $\mathcal{M}_c$ direction is unperturbed by the prior change);
  \item $\Delta q = \manifest{T3: delta q}$;
  \item $\Delta \chieff = \manifest{T3: delta chi_eff}$;
  \item $\Delta (\Lambdatilde/100) = \manifest{T3: delta Lambda}$;
  \item $\Delta \ln \Dl = \manifest{T3: delta DL}$.
\end{itemize}
Restricted to the four-dimensional $\mathcal{M}_c$-orthogonal subspace,
the alignment angle between $\Delta x$ and $\vd$ is
$\manifest{T3: alignment angle}$ degrees, with bootstrap CI
$[\manifest{T3: alignment angle lo}$, $\manifest{T3: alignment angle hi}]$.
Under the random-direction null, this yields
$p = \manifest{T3: alignment p-value}$.
The information-fraction decomposition
$\lambda_i(v_i^{\top}\Delta x)^2 / \rhosq$ at full band
[\manifest{T5: info fraction per eigendirection}] confirms
that $\manifest{T5: smallest-lambda fraction}$ of the discriminating
power is concentrated on $\vd$, as predicted by the lensing argument.

Cross-approximant persistence: the alignment angle under the
TaylorF2 release pair is $\manifest{T3: alignment angle (TaylorF2)}$,
consistent with the IMRPhenom result.

\subsection{GW150914 contrast (Test C)}

Re-analysis of GW150914 posterior samples yields
$q = \manifest{T4: GW150914 q}$ (median and 90\% CI),
$\chieff = \manifest{T4: GW150914 chi_eff}$, correcting the
previously quoted $q = 0.99 \pm 0.03$ from memory.  The Fisher
condition number for GW150914 (4D parameter space, no tidal sector)
at matched cutoffs is $\manifest{T4: condition number ratio}$ times
smaller than GW170817's, predicting a prior-induced displacement
$\manifest{T4: predicted displacement ratio}$ times smaller.
The actual displacement between alternative GW150914 analyses is
$\manifest{T4: GW150914 displacement}$ along $\vd$, consistent with
this contrast.


\section{Dynamic Test: Fisher Information Sweep and Commutator Growth}
\label{sec:dynamic}

\subsection{The prediction}

If the ($\rho$, $\Pi$) formalism describes the evolution of the
posterior under accumulating data, then the commutator
$\|[\Gamma, \Pi]\|$ should grow monotonically with frequency cutoff
$f$, driven by the PN hierarchy.  At low frequencies the Fisher matrix
is dominated by $\Mchirp$ alone; the projector $\Pi$ onto this direction
is nearly aligned with $\Gamma$, so $C(f)$ is small.  At higher
frequencies, higher PN orders become resolvable, adding structure to
the likelihood surface in $(q, \chieff, \Lambdatilde)$; $\Gamma$ tilts,
and $C(f)$ grows.  The growth law is derivable from PN power counting
and can be compared against the numerical sweep.

\subsection{Method}

We use the TaylorF2 post-Newtonian waveform approximant (3.5PN
point-particle $+$ leading-order tidal terms), implemented in JAX for
machine-precision derivatives via \texttt{jax.jacfwd} (eliminating the
step-size uncertainty of finite differences).  The parameter vector in
declared coordinates is
\begin{equation}
x = (\ln \Mchirp,\, q,\, \chieff,\, \Lambdatilde/100,\, \ln \Dl).
\label{eq:param}
\end{equation}
The Fisher matrix is computed on a grid of $N=500$ log-spaced frequency
bins over $[20, 2048]$~Hz, with cumulative trapezoidal integration.
The cutoff sweep is evaluated at 18 manuscript cutoffs plus a dense
200-point grid for slope fitting.  All $C(f)$ values are computed at
three expansion points (lowSpin mean, highSpin mean, midpoint); the
spread is reported as a systematic band.

\subsection{Results: commutator growth}

For GW170817-like parameters, $C(f)$ grows monotonically from
$\manifest{T6: C at 22 Hz}$ at 22~Hz to
$\manifest{T6: C at 500 Hz}$ at 500~Hz, a growth factor of
$\manifest{T6: growth factor}$.

\begin{table*}[tp]
\centering
\renewcommand{\arraystretch}{1.2}
\caption{Commutator growth across the chirp window.  All values in
declared coordinates (\eqnref{coords}) at the midpoint expansion point;
systematic band from ridge-uniformity check in brackets.}
\label{tab:commutator_growth}
\begin{tabular}{lccc}
\toprule
$f_{\text{max}}$ (Hz) & $\tau$ to merger (s) & $C(f)$ & Growth factor \\
\midrule
22  & 126.0 & \manifest{T6: C at 22 Hz}   & $1.00\times$ \\
35  &  36.5 & \manifest{T6: C at 35 Hz}   & \manifest{T6: growth to 35 Hz} \\
75  &   4.8 & \manifest{T6: C at 75 Hz}   & \manifest{T6: growth to 75 Hz} \\
200 &   0.4 & \manifest{T6: C at 200 Hz}  & \manifest{T6: growth to 200 Hz} \\
500 &  $\sim0$ & \manifest{T6: C at 500 Hz} & \manifest{T6: growth factor} \\
\bottomrule
\end{tabular}
\end{table*}

\subsection{Robustness battery}

The growth factor under two alternative coordinate conventions
(ALT-A: symmetric mass ratio; ALT-B: prior-range whitening) and across
three expansion points yields:
$\manifest{T6: growth factor} \pm \manifest{T6: robustness envelope}$.
A cross-check using IMRPhenomD\_NRTidal at five cutoffs gives agreement
within $\manifest{T6: approximant agreement}$.

\subsection{Derived growth law}
\label{sec:growthlaw}

The commutator growth $C(f)$ follows from PN power counting.  In
TaylorF2, the phase is a sum $\psi = \sum_k \psi_k f^{(k-5)/3}$;
each Fisher element is an integral of
$\int f^{-7/3} f^{(j+k-10)/3} S_n(f)^{-1} df$.  Under a power-law
approximation to the PSD over the band, the cumulative eigenvalues
become explicit power laws in $f$, and the off-diagonal block
connecting $\ln \Mchirp$ to $(q, \chieff, \Lambdatilde)$ grows with
extra powers of $(\pi M f)^{2/3}$ per PN order.  The misalignment
$C(f)$ is a ratio of these cumulative power laws; when the leading
exponents nearly cancel---as the PN hierarchy enforces in the band
where multiple orders contribute comparably---the ratio grows slowly,
with a logarithmic form $C(f) \propto \ln f$ to leading order.

The predicted slope $\mathrm{d}\ln C / \mathrm{d}\ln f$ from the
analytic derivation, averaged over $[22, 500]$~Hz, is
$\manifest{T7: predicted slope}$.  The numerical sweep gives
$\manifest{T7: measured slope}$, with piecewise agreement of
$\manifest{T7: slope agreement per band}$ in the
$[22,75]$, $[75,200]$, and $[200,500]$~Hz bands.  All bands agree
within the pre-registered tolerance of 25\%.

\subsection{PSD-insensitivity and universality}
\label{sec:psd}

We repeat the commutator analysis under four PSDs: ZDHP (analytic),
Early aLIGO (analytic), and measured O2-era H1 and L1 PSDs (GWOSC
strain-derived, spectral lines retained).  The commutator integral
(the accumulated bias formula integrand) shows a maximum variation
of $\manifest{T8: dual-PSD ratio}$ across all four PSDs:
\begin{equation}
\frac{\int \|[\Gamma, \Pi]\|_F\,df\;|\;_{\text{PSD}}}
{\int \|[\Gamma, \Pi]\|_F\,df\;|\;_{\text{ZDHP}}}
= \manifest{T8: dual-PSD ratio}.
\end{equation}
The fractional change in $C(f)$ is at the $\sim1\%$ level for all
cutoffs; a line-artifact scan (checking for anomalous steps at
known spectral-line frequencies) shows no excursions above
$\manifest{T8: max line artifact}$ times the local median step.

This $\sim1\%$ PSD insensitivity is not an unexplained empirical
fact; it follows from the derived growth law.  Changing $S_n(f)$
reweights the frequency measure, effectively shifting the PSD-weighted
mean frequency $f_{\text{eff}}$ by $\delta\ln f_{\text{eff}}$.
The induced change in $C$ is
$\delta C / C \approx (\mathrm{d}\ln C / \mathrm{d}\ln f) \cdot
\delta\ln f_{\text{eff}}$, and with the small slope from
\secref{growthlaw}, the predicted $\sim1\%$ variation matches the
computed value.

Because the effect is bounded at $\sim1\%$---far below inter-detector
statistical scatter---the multi-detector test proposed in earlier
drafts of this work is unperformable and is here withdrawn.  Its
replacement, the in-software ordering experiment, is described in
\secref{ordering}.


\section{Forward Model for the Time-Resolved Signature}
\label{sec:forward}

\subsection{The non-Gaussianity index}
\label{sec:ngauss}

The Fisher approximation under-estimates the posterior's spread along
$\vd$ at all cutoffs.  We quantify this via $N(f)$
(\eqnref{ngauss}), computed using posterior samples: at 500~Hz,
$N(500\,\text{Hz}) = \manifest{T9: N at 500 Hz}$,
confirming the Laplace approximation is off by a factor of
$\sim\!\sqrt{N}$ in width along $\vd$ even at full bandwidth.
This is the regime where the sequential-Gaussian bias of Object~2
(\secref{exactness}) accrues.

\subsection{Derived weight trajectory}
\label{sec:weights}

Under exact inference, the relative weight of two locations on the
ridge is not a free function---it is slaved to the accumulated
likelihood ratio.  For Gaussian noise with the truth at $\theta_1$,
the log-odds between $\theta_2$ and $\theta_1$ given data up to
cutoff $f$ is a random variable with
\begin{equation}
\mathbb{E}[\Delta\ln L(f)] = -\tfrac12\,\rhosq(f),\qquad
\Var[\Delta\ln L(f)] = \rhosq(f),
\label{eq:logodds}
\end{equation}
where $\rhosq(f) \equiv \Delta x^{\top}\Gamma(f)\Delta x$ is the
discriminating SNR$^2$ of \eqnref{rhosq}.  The weight ratio is
\begin{equation}
\frac{w_2(f)}{w_1(f)} \sim \exp\!\bigl(-\tfrac12 \rhosq(f)
+ \rhosq(f)^{1/2}\,z\bigr),\quad z\sim\mathcal{N}(0,1).
\label{eq:weight_derived}
\end{equation}
Both ridge regions are supported while $\rhosq \lesssim 1$; the
disfavoured region is exponentially suppressed once $\rhosq \gtrsim
\text{few}$.

We compute $\rhosq(f)$ exactly (via mismatch between the two
prior-conditioned mean waveforms) and via the quadratic approximation;
the ratio is $\manifest{T5: quadratic-to-exact ratio}$---we use
the exact form for all predictions.

\subsection{Peak-drop prediction}
\label{sec:peakdrop}

The within-mode width $\sigma^2(f) = (\Gamma(f)^{-1})_{\vd\vd}$ shrinks
monotonically with $f$, while $w_1 w_2$ holds roughly constant at
$\sim\!0.25$ for $\rhosq \lesssim 2$ and collapses exponentially
thereafter.  The between-mode variance along $\vd$ is
$B(f) = w_1(f) w_2(f)\,(\vd^{\top}\Delta x)^2$.  The ratio
$B(f)/\sigma^2(f)$ therefore rises (driven by Fisher narrowing) and
then drops (driven by weight resolution)---a peak-drop whose peak
frequency $f_*$ is given by $\rhosq(f_*) \approx 2$.

Solving $\rhosq(f_*) = 2$ yields
$f_* = \manifest{T9: predicted f*}$ Hz,
with a sensitivity band $[\manifest{T9: f* lo}$,
$\manifest{T9: f* hi}]$ from thresholds $\rhosq \in [1,4]$.
The peak magnitude is
$\manifest{T9: peak magnitude}$, and the noise-induced scatter
(drawn from the $z$ distribution in \eqnref{weight_derived}) gives an
event-to-event scatter of
$\pm\manifest{T9: f* scatter}$ Hz---a falsifiable population
prediction.

This replaces the earlier imposed-weight simulation (which built in
its conclusion by setting $w(f)$ by hand).  The earlier fixed-weight
run is recast as the non-Gaussianity index $N(f)$ above, which is
a measured diagnostic rather than a simulation result.

\begin{figure}[h]
\centering
\includegraphics[width=\columnwidth]{forward_model.png}
\caption{Forward model for the time-resolved signature.  Top: derived
weight ratio $w_2/w_1(f)$ from accumulated likelihood ratio.
Middle: between-mode variance $B(f)$, within-mode variance
$\sigma^2(f)$, and their ratio.  Bottom: $\rhosq(f)$ with the
$f_*$ threshold.  The predicted peak frequency $f_*$ is marked.
The shaded band shows the noise-induced scatter from 10$^4$ draws.}
\label{fig:forward}
\end{figure}


\section{Open Problems and the Decisive Test}
\label{sec:open}

\subsection{The ordering experiment}
\label{sec:ordering}

The most direct test of the ADF bias prediction (Object~2) is an
in-software ordering experiment.  For a single synthetic injection at
GW170817-midpoint parameters (ZDHP PSD, $f \in [22, 500]$~Hz split
into $B=18$ bands), we compare four inference arms:

\begin{enumerate}
  \item[(a)] \textbf{batch} --- nested sampling on the full-band
  likelihood (ground truth);
  \item[(b)] \textbf{ADF-forward} --- single-pass sequential Laplace
  in ascending frequency order;
  \item[(c)] \textbf{ADF-reverse} --- descending order;
  \item[(d)] \textbf{ADF-shuffled} --- 20 random permutations.
\end{enumerate}

Exact inference predicts (a) = (b) = (c) = (d) up to sampling error.
Our ADF bias formula predicts the offset $\Delta\theta^{(b,c,d)} -
\Delta\theta^{(a)}$ quantitatively, projected onto $\vd$, scaling with
$\int \|[\Gamma, \Pi]\|_F\,df$ along each ordering.  The pre-registered
success criterion is a Spearman rank correlation $> 0.7$ between
predicted and measured $\vd$-offsets across the 23 ADF runs.

This experiment replaces the previously proposed multi-detector
comparison, which is unperformable for three independent reasons:
(i) the $\sim1\%$ effect size (from the dual-PSD analysis) is far
below inter-detector statistical scatter; (ii) GW170817 parameter
estimation is coherent network inference, so independent per-detector
posteriors of comparable quality do not exist; and (iii) the PSD
robustness check used idealized design curves, not measured August 2017
spectral shapes with lines---so even the robustness claim was tested
against smooth approximations.  The ordering experiment requires no
astrophysical confounds and settles the question on a single laptop.

\subsection{Population survey of BNS mergers}

With O4 and O5 data, the population of detected BNS mergers will grow
to tens of events, allowing a population-level test.  The invariant
comparators from \secref{coordinates}---the likelihood cost
$\Delta x^{\top}\Gamma\Delta x$ and the $\vd$ information fraction---can
be compared across events, unlike the unit-dependent angles in the
earlier draft.  The $f_*$ prediction from \secref{peakdrop} yields an
event-to-event scatter distribution testable with $\sim10$ events.

\subsection{Summary of evidence}

Table~\ref{tab:evidence} replaces the $\kappa$ confidence scores of
earlier drafts with a typed, falsifiability-driven summary.

\begin{table*}[tp]
\centering
\renewcommand{\arraystretch}{1.3}
\caption{Summary of evidence.  Epistemic type: M = measurement on data,
C = deterministic model computation, I = illustration / forward model.
Falsification column states what observation would refute the claim.}
\label{tab:evidence}
\begin{tabularx}{\textwidth}{@{}l c c X@{}}
\toprule
Test & Type & Primary result & What would falsify \\
\midrule
Ridge orientation (Test A) & M &
$\manifest{T2: ridge alignment angle}$ &
Misalignment $>20^\circ$ at 95\% CI \\
Prior-displacement alignment (Test B) & M &
$p = \manifest{T3: alignment p-value}$ &
$p > 0.05$ or bootstrap CI includes random direction \\
GW150914 contrast (Test C) & M &
Displacement ratio $\manifest{T4: predicted displacement ratio}$ &
No contrast between events \\
$C(f)$ growth factor & C &
$\manifest{T6: growth factor}\pm\manifest{T6: robustness envelope}$ &
Envelope includes $1.0\times$ \\
Growth-law slope & C &
Predicted-vs-measured within 25\% per band &
Any band outside tolerance \\
PSD insensitivity & C &
Ratio $\manifest{T8: dual-PSD ratio}$ &
Anomalous step at line frequencies \\
Peak-drop prediction $f_*$ & I &
$f_* = \manifest{T9: predicted f*}$ Hz &
No peak in time-resolved analysis \\
\bottomrule
\end{tabularx}
\end{table*}

\subsection{Extension to BBH and EMRI systems (low priority)}

The formalism is general; BBH systems with precession or higher modes
may show commutator signatures in spin-precession subspaces, and
EMRI signals in the LISA band provide the ultimate test of sequential
inference geometry over $\sim10^5$ orbits.


\section{Interpretation}
\label{sec:interpretation}

\subsection{The layer of description}

The central result of this work is that the ($\rho$, $\Pi$) dynamical
system---a general model of commutator-modulated relaxation toward a
fixed point---describes the \emph{geometry of Bayesian inference} in a
gravitational-wave detector, not the binary system itself.  This is an
epistemological claim, not an ontological one: the ($\rho$, $\Pi$)
formalism does not replace General Relativity nor describe new physics
in the binary's orbital dynamics.  It describes how the posterior
distribution reorganises itself as data accumulates, how its geometry
evolves on the manifold of orbital parameters, and how this evolution
has its own measurable dynamics---including a quantitative prediction
for the bias of sequential-Gaussian pipelines that \emph{are} deployed
in low-latency analyses today.

\subsection{Why the timescales coincide}

The deepest structural result of the Fisher sweep is not a number but a
structural fact: the commutator grows on the same timescale as the
orbital chirp.  This is a \emph{constraint}: a necessary consequence of
the fact that the same strain data drives both.  The orbital chirp
($df/dt \propto f^{11/3}$) is a prediction of GR; the commutator growth
is a prediction of information geometry.  The two are locked because
each frequency maps to a time-to-merger via the GR chirp, and the
Fisher matrix at frequency $f$ knows about the waveform up to
$f$---which is the waveform at time $\tau(f)$ before merger.  No
hypothesis about modified gravity is needed or supported.

\subsection{The $\nu$ exponent}

Earlier drafts compared a claimed $\nu = 9$ exponent from the
($\rho$, $\Pi$) linearized dynamics to GR's $11/3$.  This comparison
is not meaningful: the inference system has no autonomous clock, its
independent variable is $f$ which advances because the binary chirps,
and asking for the ($\rho$, $\Pi$) system's $df/dt$ confuses the
epistemological and ontological readings.  The ($\rho$, $\Pi$) formalism
makes no claims about waveform dynamics; there is no exponent to
compare against GR.  This section is removed entirely.  (The power
counting of \secref{growthlaw} provides the well-defined rate law that
the earlier $\nu$ was attempting to supply.)

\subsection{Path-dependence of the posterior}

The exact posterior $\rho_{\fmax}$ is path-independent
(Proposition~\ref{prop:exact}).  What is path-dependent is (i) the
trajectory $\rho_f$, whose geometry $C(f)$ describes, and (ii) the
error of any single-pass Gaussian approximation, which the ADF bias
formula quantifies.  The PSD enters only because different PSDs define
different inference problems (different weighting of the likelihood
factors), not because they create different paths through the same
problem.  This distinction, unmarked in the earlier draft, is enforced
throughout this revision.

\subsection{The bridge as a methodological tool}

Beyond its specific results, the ($\rho$, $\Pi$) bridge demonstrates
how a dynamical system from one context (the S3 agent framework) can be
reinterpreted as a description of Bayesian inference in another domain.
The key step is the identification of the posterior $\rho$ and the
likelihood structure $\Pi$ with the dynamical pair.  Any sequential
Bayesian inference problem---where data arrives incrementally and the
posterior evolves---may admit a ($\rho$, $\Pi$) description, with the
commutator $[\Gamma, \Pi]$ measuring the cross-coupling between the
best-constrained direction and the remaining parameter space in that
domain.

\section*{Acknowledgments}

We thank the LIGO Scientific Collaboration and the Virgo Collaboration
for making posterior samples publicly available through the
Gravitational Wave Open Science Center (GWOSC), and Lark Laflamme for
editorial review that substantially improved the manuscript.  This work
was supported by RavenNest Science.


\begin{thebibliography}{99}

\bibitem{abbott2017gw170817}
B.P.\ Abbott et al.\ (LIGO Scientific Collaboration and Virgo Collaboration),
``GW170817: Observation of Gravitational Waves from a Binary Neutron Star Inspiral,''
Phys.\ Rev.\ Lett.\ \textbf{119}, 161101 (2017).

\bibitem{abbott2016gw150914}
B.P.\ Abbott et al.\ (LIGO Scientific Collaboration and Virgo Collaboration),
``Observation of Gravitational Waves from a Binary Black Hole Merger,''
Phys.\ Rev.\ Lett.\ \textbf{116}, 061102 (2016).

\bibitem{abbott2019gw170817_params}
B.P.\ Abbott et al.\ (LIGO Scientific Collaboration and Virgo Collaboration),
``Properties of the binary neutron star merger GW170817,''
Phys.\ Rev.\ X \textbf{9}, 011001 (2019).

\bibitem{cutler1994fisher}
C.\ Cutler and E.E.\ Flanagan,
``Gravitational waves from merging compact binaries: How accurately can
one extract the binary's parameters from the inspiral waveform?''
Phys.\ Rev.\ D \textbf{49}, 2658 (1994).

\bibitem{vallisneri2008fisher}
M.\ Vallisneri,
``Use and abuse of the Fisher information matrix in the assessment of
gravitational-wave parameter-estimation prospects,''
Phys.\ Rev.\ D \textbf{77}, 042001 (2008).

\bibitem{poisson2014gravity}
E.\ Poisson and C.M.\ Will,
\emph{Gravity: Newtonian, Post-Newtonian, Relativistic}
(Cambridge University Press, 2014).

\bibitem{buonanno2009gw}
A.\ Buonanno and B.S.\ Sathyaprakash,
``Sources of Gravitational Waves: Theory and Observations,''
arXiv:0906.5312 (2009).

\bibitem{maggiore2007gw}
M.\ Maggiore,
\emph{Gravitational Waves: Volume 1: Theory and Experiments}
(Oxford University Press, 2007).

\bibitem{flanagan1998fisher}
E.E.\ Flanagan and S.A.\ Hughes,
``Measuring gravitational waves from binary black hole coalescences:
I.\ Signal-to-noise for inspiral, merger, and ringdown,''
Phys.\ Rev.\ D \textbf{57}, 4535 (1998).

\bibitem{chatziioannou2020bns}
K.\ Chatziioannou,
``Uncertainty in the neutron star equation of state from gravitational
waves,''
arXiv:2006.03168 (2020).

\bibitem{minka2001ep}
T.P.\ Minka,
``Expectation propagation for approximate Bayesian inference,''
in \emph{Proc.\ 17th Conf.\ on Uncertainty in Artificial Intelligence}
(UAI), 362 (2001).

\bibitem{opper1998online}
M.\ Opper,
``A Bayesian approach to on-line learning,''
in \emph{On-Line Learning in Neural Networks},
edited by D.\ Saad (Cambridge University Press, 1998), p.\ 363.

\end{thebibliography}


\appendix

\section{Eigenframe-leakage derivation of the ADF bias}
\label{sec:appendix}

Let the exact batch posterior at full band have mode $\hat\theta$ and
precision $\Gamma = \Gamma(f_{\text{max}})$.  Consider a single-pass
sequential Laplace (ADF) recursion that processes frequency bands in
order $B_1, \ldots, B_N$.  After band $k$ the filter maintains a
Gaussian $q_k(\theta) = \mathcal{N}(m_k, P_k)$, where $P_0 = \Gamma_0^{-1}$
for some initial prior curvature and $m_0 = \theta_{\text{prior}}$.

Upon receiving band $B_{k+1}$ with incremental Fisher matrix
$\Delta\Gamma_{k+1}$ (computed at the true parameters $\theta_0$),
the filter linearises $h$ about $m_k$ and performs a conjugate update:
\begin{equation}
P_{k+1}^{-1} = P_k^{-1} + \Delta\Gamma_{k+1}, \qquad
m_{k+1} = m_k + P_{k+1}\,\bigl[\nabla\ell(m_k)_{\text{new}}\bigr],
\label{eq:adf_update}
\end{equation}
where $\nabla\ell(m_k)_{\text{new}}$ is the score of the new band
evaluated at $m_k$.  The exact score at the batch mode vanishes:
$\sum_k \nabla\ell_k(\hat\theta) = 0$.  The bias $\delta_k = m_k - \hat\theta$
evolves as
\begin{equation}
\delta_{k+1} = \delta_k + P_{k+1}\,[\nabla\ell_{k+1}(m_k)
- \nabla\ell_{k+1}(\hat\theta)] + P_{k+1}\nabla\ell_{k+1}(\hat\theta)
- P_k^{-1}\delta_k .
\label{eq:bias_recurrence}
\end{equation}

To first order in $\delta_k$, the term in square brackets is
$-\Delta\Gamma_{k+1}\,\delta_k$, and the last two terms cancel at
the linearised level when $P_{k+1}^{-1}\delta_k = (P_k^{-1} +
\Delta\Gamma_{k+1})\delta_k = P_k^{-1}\delta_k + \Delta\Gamma_{k+1}\delta_k$.
The recurrence simplifies to
\begin{equation}
\delta_{k+1} \approx P_{k+1}P_k^{-1}\delta_k
+ P_{k+1}\,\nabla\ell_{k+1}(\hat\theta),
\label{eq:bias_simple}
\end{equation}
which propagates past errors forward.  Crucially, if every increment
$\Delta\Gamma_{k+1}$ commuted with the cumulative $\Gamma_k$ (i.e.,
shared its eigenframe), then $P_{k+1}P_k^{-1} = (I + \Gamma_k^{-1}
\Delta\Gamma_{k+1})^{-1}$ would be diagonal in the common eigenbasis,
and the displacement $\delta_k$ along a well-constrained direction
would not leak into the degenerate subspace.  The off-diagonal part of
$\Delta\Gamma_{k+1}$ in $\Gamma_k$'s eigenframe---measured by
$\|[\Gamma_k, \Delta\Gamma_{k+1}]\|$---is the leakage operator.

Integrating over the continuous frequency path (the continuum limit of
the band recursion) and projecting the accumulated bias onto the
degeneracy direction $\vd$ gives
\begin{equation}
\Delta\theta_{\text{sys}} \approx
\Gamma(f_{\text{max}})^{-1}\,
\int_{f_{\text{min}}}^{f_{\text{max}}}
[\Gamma(f'), \Pi(f')]\,df' \cdot \vd .
\label{eq:sys_offset}
\end{equation}

This is the ADF bias formula.  The commutator controls the rate at
which sequential approximation error is rotated into the degeneracy
direction; the factor $\Gamma^{-1}$ amplifies it along $\vd$ where
the eigenvalue is smallest; and the integral accumulates the leakage
along the frequency path.  The expression also supplies the
commutation condition for vanishing bias: if $[\Gamma, \Pi] = 0$ for
all $f$ (i.e., the eigenframe does not rotate), the bias is zero.

\section{Conditioning diagnostics}
\label{sec:appendix_conditioning}

All eigendecompositions are performed in the declared coordinates
(\eqnref{coords}).  The condition number $\kappa(\Gamma(f))$ across
the sweep is $\manifest{T11: condition number range}$.  The smallest
spectral gap $\lambda_4(f) - \lambda_5(f)$ is
$\manifest{T11: min spectral gap}$; eigenvector $\vd$ stability is
verified by perturbing $\Gamma$ at the level of the numerical
precision of the JAX derivatives and measuring the resulting angle
change ($\manifest{T11: vd perturbation angle}$).

A finite-difference cross-check (using lalsuite at 5 step sizes)
against the JAX autodiff result shows that the vd angle error from
finite differences ranges from $\manifest{T11: FD angle error min}$
to $\manifest{T11: FD angle error max}$ degrees across step sizes,
confirming the necessity of machine-precision derivatives for the
degenerate eigenvalue problem.

\end{document}